\begin{proof}
\begin{enumerate}[label={$\langle1\rangle$\arabic*.},leftmargin=*]
\item 
\noindent\textbf{Step 1.} The section is nonzero and maps to zero in the quotient
$\widetilde\Omega^1_S=\Omega^1_S/\mathscr T_S$, so it
lies in the left-hand side of the section-kernel isomorphism of \cref{cdp:st:A.1.1}.  Proposition~\ref{cdp:st:A.1.1}
identifies that group with $\ker\delta_A$, which is therefore nonzero.

\end{enumerate}
\end{proof}

\begin{proof}
\begin{enumerate}[label={$\langle1\rangle$\arabic*.},leftmargin=*]
\item 
Let $\Delta$ be a sufficiently small disc with coordinate $t$ and put
\[
\mathcal C=
\{([x:y:z],t)\in\mathbb P^2\times\Delta:
y^2z=x^2(x+z)+tz^3\}.
\tag{W.8}\hypertarget{tag-03I1}{}\label{wbII:eq:W:5.8}
\]
This is proper over $\Delta$.  Its total space is smooth.  On $z\ne0$, the
partial derivative of the defining equation with respect to $t$ is
$-z^3\ne0$.  On $z=0$, the equation forces $x=0$; then $y\ne0$ and the
$z$-derivative is $y^2\ne0$.  The central fibre $C_0$ is the irreducible
nodal cubic $y^2z=x^2(x+z)$, with node $[0:0:1]$.

Set
\[
\mathcal X=\mathcal C\times\mathbb P^1,\qquad
S=C_0\times\mathbb P^1,
\]
and let $g:\mathcal X\to\Delta$ be projection.  Then $\mathcal X$ is a
smooth threefold, $g$ is proper, and $S$ is an irreducible reduced
normal-crossings surface.  Its normalization is
$\widetilde S\simeq\mathbb P^1\times\mathbb P^1$.

Let $p,q$ be the two inverse images of the node in the normalization of
$C_0$.  Fix $\lambda\in\mathbb C^*$ which is not a root of unity (for
example, $\lambda=2$).  Glue the fibres of
the trivial line bundle on the normalization at $p$ and $q$ by the scalar
$\lambda$.  This defines a line bundle
$A_\lambda\in\operatorname{Pic}^0(C_0)$ whose pullback is trivial.  It has
infinite order: its $n$th tensor power has gluing scalar $\lambda^n$, which
is never one.  It is not trivial: a global section pulls back to a constant
$c$ satisfying
$c=\lambda c$, so
\[
H^0(C_0,A_\lambda)=0.
\tag{W.9}\hypertarget{tag-03I2}{}\label{wbII:eq:W:5.9}
\]
Put $A=\operatorname{pr}_1^*A_\lambda$ on $S$.  Then
$\eta^*A\simeq\mathcal O_{\widetilde S}$ and $H^0(S,A)=0$.
Because $S$ is the full fibre,
$N^*_{S/\mathcal X}\simeq\mathcal O_S$, so the conormal-twist global
sections vanish.

The natural map
$\widetilde\Omega^1_S\to\eta_*\Omega^1_{\widetilde S}$ is injective:
its kernel is supported where normalization is not an isomorphism, while
the source is torsion-free by definition.  Projection formula and
$\eta^*A\simeq\mathcal O_{\widetilde S}$ therefore give
\[
H^0(S,\widetilde\Omega^1_S\otimes A)
\hookrightarrow
H^0(\widetilde S,\Omega^1_{\widetilde S})=0,
\tag{W.10}\hypertarget{tag-03EA}{}\label{wbII:eq:W:5.10}
\]
the last equality following from
$\widetilde S=\mathbb P^1\times\mathbb P^1$.
Nevertheless \cref{wbII:prop:conductor} produces
\[
0\ne s\in H^0(S,\Omega^1_{\mathcal X}|_S\otimes A).
\]
Thus both vanishings invoked in the printed reduction hold while the desired
ambient vanishing fails.  This calculation chooses $A$ directly on $S$ and disproves the
claimed \emph{formal implication between the three sheaf-cohomology groups}
for a non-torsion line bundle. The explicit extending choice of $A$ is
constructed in the next subsection, which also treats the total-space
line-bundle condition on p.~692.  The family \eqref{wbII:eq:W:5.8} also does not satisfy all hypotheses of CDP
Theorem~2.2, so this is not a counterexample to that theorem.

\end{enumerate}
\end{proof}

\begin{proof}
\begin{enumerate}[label={$\langle1\rangle$\arabic*.},leftmargin=*]
\item 
Properness follows because $\mathcal C$ is closed in the projective
projection $\PP^2\times\Delta\to\Delta$.  Write the defining equation as
\[
 F=y^2z-x^2(x+z)-tz^3.
\]
On the chart $z\ne0$, $\partial F/\partial t=-z^3\ne0$, so $\mathcal C$
is smooth there.  On $z=0$, the equation forces $x=0$ and $y\ne0$; at such
points $\partial F/\partial z=y^2\ne0$.  Hence $\mathcal C$, and therefore
$\mathcal X$, is smooth.  The only finite discriminant values of the fibres
are $t=0$ and $t=-4/27$; after shrinking $\Delta$ to $|t|<4/27$, every
noncentral fibre is smooth.

The central curve is
\(C_0:\ y^2z=x^2(x+z)\).  Its only singular point is
\(n=[0:0:1]\): on $z=1$ the equations $F=F_x=F_y=0$ give $(x,y)=(0,0)$,
and the point at infinity $[0:1:0]$ is smooth because $F_z=1$ there
(Euler's identity also gives the remaining projective singularity equation).
The parametrization
\[
 \nu:\PP^1\longrightarrow C_0,\qquad
 [u:v]\longmapsto
 [(u^2-v^2)v:\,u(u^2-v^2):\,v^3]
 \tag{W.160}\hypertarget{tag-03DX}{}\label{wbII:eq:CN:4.3}
\]
maps the two points $p=[1:1]$, $q=[-1:1]$ to $n$.  On the chart $v\ne0$, set
$r=u/v$; after scaling, its image is
$[r^2-1:r(r^2-1):1]$.  Conversely, for an affine point of $C_0$ with
$x\ne0$, the equation gives $r=y/x$ and $x=r^2-1$, so it is hit by
$[r:1]$.  The only affine point with $x=0$ is the node $(0,0)$, hit by
$[1:1]$ and $[-1:1]$, and the point at infinity $[0:1:0]$ is hit by
$[1:0]$.  Thus $\nu(\PP^1)=C_0$; as a nonconstant morphism to a projective
curve it is finite and generically one-to-one, hence the normalization.  In
particular $C_0$ is irreducible and rational.  The product normalization
is $\eta=\nu\times\id_{\PP^1}:\PP^1\times\PP^1\to S$.

At $n$ the affine equation is
\(t=y^2-x^2-x^3\).  On a sufficiently small neighbourhood choose the
holomorphic square root $h(x)=\sqrt{1+x}$ and set
\(u=y-xh(x)\), \(v=y+xh(x)\).  Then $uv=t$, and the map
$(x,y)\mapsto(u,v)$ has invertible $2\times2$ Jacobian at $(0,0)$, so $u,v$
are local coordinates on the surface $\mathcal C$.  If $s$ is a local
coordinate on $\PP^1$, then $(u,v,s)$ are coordinates on $\mathcal X$ and
$t=uv$.  Hence the pair has the asserted normal-crossing (self-crossing)
local model, with singular locus $D$ as stated.

\end{enumerate}
\end{proof}
